% Generated by scripts/assemble.py; corrections belong in source/corrections.json.
\examhead{93}{南京师范大学 2026 年数学分析试卷}{92}
% p093-b01; source PDF page 93
\par\Needspace{5\baselineskip}
\blockanchor{p093-b01}
\majorquestion{一}计算题. 每题6分, 共30分.\par
% p093-b02; source PDF page 93
\blockanchor{p093-b02}
\question{1}求极限  $\lim_{n\to\infty}\left[\frac{1}{(n+1)(1+\frac{1}{n})}+\frac{1}{(n+\frac{1}{2})(1+\frac{2}{n})}+\cdots+\frac{1}{(n+\frac{1}{n})(1+\frac{n}{n})}\right]$\par
% p093-b03; source PDF page 93
\blockanchor{p093-b03}
\question{2}求积分 $\int_0^\pi \frac{x\sin x}{4\cos^2x + 3\sin^2x}\mathrm{d}x.$\par
% p093-b04; source PDF page 93
\blockanchor{p093-b04}
\question{3}求重积分 $\iint_{D} f(x, y) \, \mathrm{d}x \, \mathrm{d}y$, 其中 $f(x, y) = \begin{cases} y, & y \leq e^x, \\ x, & y > e^x. \end{cases}$$D: [0,1] \times [0,e]$.\par
% p093-b05; source PDF page 93
\blockanchor{p093-b05}
\question{4}求由 $y^{2} = 2x, y^{2} = 4x, y = x, y = 2x, z = 0, z = xy$ 围成区域的体积.\par
% p093-b06; source PDF page 93
\blockanchor{p093-b06}
\question{5}设 $f(x, y) = x^y \ln(xy)$，求 $f(x, y)$ 在 (1, 1) 处的二阶泰勒公式.\par
% p093-b07; source PDF page 93
\par\Needspace{5\baselineskip}
\blockanchor{p093-b07}
\majorquestion{二}(15 分) 设 $f(x)$ 在 $[0,2]$ 上恒正连续, $a_{n} = \int_{0}^{2} x^{n} f(x) \, \mathrm{d}x$, 求幂级数 $\sum_{n=0}^{\infty} a_{n} x^{n}$ 的收敛域.\par
% p093-b08; source PDF page 93
\par\Needspace{5\baselineskip}
\par\Needspace{6\baselineskip}
\blockanchor{p093-b08}
\majorquestion{三}(15 分) 设\par
\nopagebreak[4]
% p093-b09; source PDF page 93
\blockanchor{p093-b09}
\displayformula{f (x, y) = \left\{ \begin{array}{l l} \frac {\sin (x y)}{y}, & y \neq 0; \\ 0, & y = 0. \end{array} \right.}
% p093-b10; source PDF page 93
\blockanchor{p093-b10}
判断 $f$ 在原点 $(0,0)$ 处是否连续，可微？方向导数是否存在？\par
% p093-b11; source PDF page 93
\par\Needspace{5\baselineskip}
\blockanchor{p093-b11}
\majorquestion{四}(15 分) 求曲线积分 $\oint_{L} \frac{x \, \mathrm{d}y - y \, \mathrm{d}x}{4x^{2} + y^{2}}$, 其中 $L: (x - 1)^{2} + y^{2} = 4$, 取逆时针方向.\par
% p093-b12; source PDF page 93
\par\Needspace{5\baselineskip}
\par\Needspace{6\baselineskip}
\blockanchor{p093-b12}
\majorquestion{五}(15 分) 求曲面积分\par
\nopagebreak[4]
% p093-b13; source PDF page 93
\blockanchor{p093-b13}
\displayformula{\iint_ {\Sigma} \frac {x \mathrm{d} y \mathrm{d} z + y \mathrm{d} z \mathrm{d} x + z \mathrm{d} x \mathrm{d} y}{\sqrt {(x ^ {2} + y ^ {2} + z ^ {2}) ^ {3}}}.}
% p093-b14; source PDF page 93
\blockanchor{p093-b14}
其中 $\Sigma$ 为 $1 - \frac{z}{9} = \frac{(x - 1)^2}{4} +\frac{(y - 2)^2}{9}$ 在 $xOy$ 平面上方的部分，取上侧.\par
% p093-b15; source PDF page 93
\par\Needspace{5\baselineskip}
\blockanchor{p093-b15}
\majorquestion{六}(15 分) 证明: $f(x)$ 在有界区间 $I$ 上一致连续的充要条件是对任意 $I$ 中的柯西列 $\{a_{n}\}$, 有 $\{f(a_{n})\}$ 也是柯西列.\par
% p093-b16; source PDF page 93
\par\Needspace{5\baselineskip}
\blockanchor{p093-b16}
\majorquestion{七}(15 分) 已知点 $P_{0}(x_{0}, y_{0}, z_{0})$ 满足方程 $e^{-xy} + 2z - e^{z} = 0$, 添加一个充分条件, 使得 $P_{0}$ 附近确定一个隐函数 $z = z(x, y)$, 并问这个隐函数有没有极大(小)值?\par
% p093-b17; source PDF page 93
\par\Needspace{5\baselineskip}
\blockanchor{p093-b17}
\majorquestion{八}(15 分) 设 $f(x)$ 在 $[-1, 1]$ 上二阶可导, 且 $f(0) = f'(0) = 0$, $f(-1) = -1$, $f(1) = 1$.\par
% p093-b18; source PDF page 93
\blockanchor{p093-b18}
(1) 证明: 存在 $\xi \in (-1,1)$, 使得 $|f''(\xi)| < 2$.\par
% p093-b19; source PDF page 93
\blockanchor{p093-b19}
(2) 证明: 存在 $\eta \in (-1,1)$, 使得 $|f''(\eta)| > 2$.\par
% p093-b20; source PDF page 93
\par\Needspace{5\baselineskip}
\blockanchor{p093-b20}
\majorquestion{九}(15 分) 设 $I(y) = \int_{0}^{+\infty} \frac{\sin^3 x}{x^y} \, \mathrm{d}x$.\par
% p093-b21; source PDF page 93
\blockanchor{p093-b21}
(1) 求 $I(y)$ 的收敛域.\par
% p093-b22; source PDF page 93
\blockanchor{p093-b22}
(2) 判断  $I(y)$  在收敛域上是否一致收敛.\par
% p093-b23; source PDF page 93
\blockanchor{p093-b23}
(3) 求 $I(3)$.\par
